% GENERATED FILE. Edit canonical lessons, then run npm run generate:books.
% canonical-source-hash: fnv1a64:512dc83104517e85
% canonical-lessons: SA-C388-vyavaharati, SA-C388-upeksate, SA-C388-saumyah, SA-C388-vrtha, SA-C388-adarsah

\chapter{Dealing with People}
\label{ch:sa-dealing-with-people}
\hlchaptermodality{388}

\begin{chapteropening}
\textbf{By the end of this chapter:} \emph{I can say how someone will be treated if they behave well or badly.}

Complete the last lesson of chapter 388: I can say how someone will be treated if they behave well or badly.
\end{chapteropening}

\section[\texorpdfstring{\sk{व्यवहरति} (vyavaharati)}{व्यवहरति (vyavaharati)}]{\sk{व्यवहरति} (vyavaharati) — deals with people; behaves}

\label{lesson:SA-C388-vyavaharati}



\begin{quote}
Before the new one: say the Sanskrit for refines, then the Sanskrit for growth.
\end{quote}

\subsection*{You'll want to know: \sk{व्यवहरति}}
\textbf{\sk{व्यवहरति}} — \emph{vyavaharati} — \textquotedblleft{}deals with people; behaves\textquotedblright{}.

\begin{quote}
\sk{यदि} \sk{सः} \sk{विनयेन} \sk{व्यवहरति}, \sk{तर्हि} \sk{जनाः} \sk{तं} \sk{सम्मानयन्ति।}
\end{quote}

\begin{quote}
\emph{yadi saḥ vinayena vyavaharati, tarhi janāḥ taṃ sammānayanti.}
\end{quote}

\begin{quote}
If he deals with people politely, they honour him.
\end{quote}

Past (he …): \sk{सः} \sk{विनयेन} \sk{व्यवहृतवान्}, \emph{saḥ vinayena vyavahṛtavān}.

\begin{cousinweb}
\textbf{\sk{व्यवहरति}} (\emph{vyavaharati}) is built on \emph{vi}, \emph{ava} and the root √hṛ, \textquotedblleft{}to take\textquotedblright{}; MW gives \emph{vy-ava-√hṛ} \textquotedblleft{}to act, to behave towards or deal with\textquotedblright{} (MW).
\end{cousinweb}

\begin{grammarlens}[title={What you've built}]
One more word for saying what happens if.
\end{grammarlens}

\subsection*{Your turn}
\begin{itemize}
\raggedright
  \item \emph{Say it:} \emph{vyavaharati}
  \item \emph{Say it:} \emph{vyavaharati}, once more
  \item \emph{Say it:} \emph{yadi saḥ vinayena vyavaharati, tarhi janāḥ taṃ sammānayanti.} — If he deals with people politely, they honour him.
  \item \emph{Recall:} say the Sanskrit for refines, then the Sanskrit for growth, then say \emph{vyavaharati} again
\end{itemize}

\begin{tcolorbox}[breakable,colback=teal!4,colframe=teal!35!black,title={Before you move on}]
What does \sk{व्यवहरति} mean? (\textquotedblleft{}Deals with people; behaves\textquotedblright{}.) Say it once more.
\end{tcolorbox}
\section[\texorpdfstring{\sk{उपेक्षते} (upekṣate)}{उपेक्षते (upekṣate)}]{\sk{उपेक्षते} (upekṣate) — neglects; overlooks}

\label{lesson:SA-C388-upeksate}



\begin{quote}
Before the new one: say the Sanskrit for growth, then the Sanskrit for deals with people.
\end{quote}

\subsection*{You'll want to know: \sk{उपेक्षते}}
\textbf{\sk{उपेक्षते}} — \emph{upekṣate} — \textquotedblleft{}neglects; overlooks\textquotedblright{}.

\begin{quote}
\sk{यदि} \sk{सः} \sk{स्वास्थ्यम्} \sk{उपेक्षते}, \sk{तर्हि} \sk{रुग्णः} \sk{भविष्यति।}
\end{quote}

\begin{quote}
\emph{yadi saḥ svāsthyam upekṣate, tarhi rugṇaḥ bhaviṣyati.}
\end{quote}

\begin{quote}
If he neglects his health, he will fall ill.
\end{quote}

Past (he …): \sk{सः} \sk{कार्यम्} \sk{उपेक्षितवान्}, \emph{saḥ kāryam upekṣitavān}.

\begin{cousinweb}
\textbf{\sk{उपेक्षते}} (\emph{upekṣate}) is built on \emph{upa} and the root √īkṣ, \textquotedblleft{}to see\textquotedblright{}; MW gives \emph{upekṣ} \textquotedblleft{}to overlook, to disregard, to neglect\textquotedblright{} (MW).
\end{cousinweb}

\begin{grammarlens}[title={What you've built}]
One more word for saying what happens if.
\end{grammarlens}

\subsection*{Your turn}
\begin{itemize}
\raggedright
  \item \emph{Say it:} \emph{upekṣate}
  \item \emph{Say it:} \emph{upekṣate}, once more
  \item \emph{Say it:} \emph{yadi saḥ svāsthyam upekṣate, tarhi rugṇaḥ bhaviṣyati.} — If he neglects his health, he will fall ill.
  \item \emph{Recall:} say the Sanskrit for growth, then the Sanskrit for deals with people, then say \emph{upekṣate} again
\end{itemize}

\begin{tcolorbox}[breakable,colback=teal!4,colframe=teal!35!black,title={Before you move on}]
What does \sk{उपेक्षते} mean? (\textquotedblleft{}Neglects; overlooks\textquotedblright{}.) Say it once more.
\end{tcolorbox}
\section[\texorpdfstring{\sk{सौम्यः} (saumyaḥ)}{सौम्यः (saumyaḥ)}]{\sk{सौम्यः} (saumyaḥ) — gentle; mild-mannered}

\label{lesson:SA-C388-saumyah}



\begin{quote}
Before the new one: say the Sanskrit for deals with people, then the Sanskrit for neglects.
\end{quote}

\subsection*{You'll want to know: \sk{सौम्यः}}
\textbf{\sk{सौम्यः}} — \emph{saumyaḥ} — \textquotedblleft{}gentle; mild-mannered\textquotedblright{}.

\begin{quote}
\sk{यदि} \sk{नायकः} \sk{सौम्यः}, \sk{तर्हि} \sk{जनाः} \sk{तम्} \sk{अनुगच्छन्ति।}
\end{quote}

\begin{quote}
\emph{yadi nāyakaḥ saumyaḥ, tarhi janāḥ tam anugacchanti.}
\end{quote}

\begin{quote}
If a leader is gentle, people follow him.
\end{quote}

\begin{cousinweb}
\textbf{\sk{सौम्यः}} (\emph{saumyaḥ}) is built on \emph{soma}, the Soma drink and the moon, with the ending \emph{-ya}; MW gives \emph{saumya} \textquotedblleft{}relating to Soma\textquotedblright{}, then \textquotedblleft{}resembling the moon\textquotedblright{}, and so \textquotedblleft{}placid, gentle, mild\textquotedblright{} (MW).
\end{cousinweb}

\begin{grammarlens}[title={What you've built}]
One more word for saying what happens if.
\end{grammarlens}

\subsection*{Your turn}
\begin{itemize}
\raggedright
  \item \emph{Say it:} \emph{saumyaḥ}
  \item \emph{Say it:} \emph{saumyaḥ}, once more
  \item \emph{Say it:} \emph{yadi nāyakaḥ saumyaḥ, tarhi janāḥ tam anugacchanti.} — If a leader is gentle, people follow him.
  \item \emph{Recall:} say the Sanskrit for deals with people, then the Sanskrit for neglects, then say \emph{saumyaḥ} again
\end{itemize}

\begin{tcolorbox}[breakable,colback=teal!4,colframe=teal!35!black,title={Before you move on}]
What does \sk{सौम्यः} mean? (\textquotedblleft{}Gentle; mild-mannered\textquotedblright{}.) Say it once more.
\end{tcolorbox}
\section[\texorpdfstring{\sk{वृथा} (vṛthā)}{वृथा (vṛthā)}]{\sk{वृथा} (vṛthā) — in vain}

\label{lesson:SA-C388-vrtha}



\begin{quote}
Before the new one: say the Sanskrit for neglects, then the Sanskrit for mild-mannered.
\end{quote}

\subsection*{You'll want to know: \sk{वृथा}}
\textbf{\sk{वृथा}} — \emph{vṛthā} — \textquotedblleft{}in vain\textquotedblright{}.

\begin{quote}
\sk{यदि} \sk{भवान्} \sk{न} \sk{शृणोति}, \sk{तर्हि} \sk{मम} \sk{उपदेशः} \sk{वृथा।}
\end{quote}

\begin{quote}
\emph{yadi bhavān na śṛṇoti, tarhi mama upadeśaḥ vṛthā.}
\end{quote}

\begin{quote}
If you do not listen, my advice is in vain.
\end{quote}

\begin{cousinweb}
\textbf{\sk{वृथा}} (\emph{vṛthā}) MW gives \emph{vṛthā} \textquotedblleft{}at will\textquotedblright{} and \textquotedblleft{}in vain, uselessly, fruitlessly\textquotedblright{}; MW gives no derivation, so its origin is uncertain (MW).
\end{cousinweb}

\begin{grammarlens}[title={What you've built}]
One more word for saying what happens if.
\end{grammarlens}

\subsection*{Your turn}
\begin{itemize}
\raggedright
  \item \emph{Say it:} \emph{vṛthā}
  \item \emph{Say it:} \emph{vṛthā}, once more
  \item \emph{Say it:} \emph{yadi bhavān na śṛṇoti, tarhi mama upadeśaḥ vṛthā.} — If you do not listen, my advice is in vain.
  \item \emph{Recall:} say the Sanskrit for neglects, then the Sanskrit for mild-mannered, then say \emph{vṛthā} again
\end{itemize}

\begin{tcolorbox}[breakable,colback=teal!4,colframe=teal!35!black,title={Before you move on}]
What does \sk{वृथा} mean? (\textquotedblleft{}In vain\textquotedblright{}.) Say it once more.
\end{tcolorbox}
\section[\texorpdfstring{\sk{आदर्शः} (ādarśaḥ)}{आदर्शः (ādarśaḥ)}]{\sk{आदर्शः} (ādarśaḥ) — an ideal; a model}

\label{lesson:SA-C388-adarsah}



\begin{quote}
Before the new one: say the Sanskrit for mild-mannered, then the Sanskrit for in vain.
\end{quote}

\subsection*{You'll want to know: \sk{आदर्शः}}
\textbf{\sk{आदर्शः}} — \emph{ādarśaḥ} — \textquotedblleft{}an ideal; a model\textquotedblright{}.

\begin{quote}
\sk{यदि} \sk{पिता} \sk{सत्यवादी}, \sk{तर्हि} \sk{सः} \sk{पुत्रस्य} \sk{आदर्शः} \sk{भवति।}
\end{quote}

\begin{quote}
\emph{yadi pitā satyavādī, tarhi saḥ putrasya ādarśaḥ bhavati.}
\end{quote}

\begin{quote}
If a father is truthful, he is a model for his son.
\end{quote}

A masculine noun.

\begin{cousinweb}
\textbf{\sk{आदर्शः}} (\emph{ādarśaḥ}) is built on \emph{ā} and the root √dṛś, \textquotedblleft{}to see\textquotedblright{}; MW gives \emph{ā-darśa} \textquotedblleft{}a mirror\textquotedblright{}, \textquotedblleft{}a copy\textquotedblright{} and \textquotedblleft{}ideal perfection\textquotedblright{} (MW).
\end{cousinweb}

\begin{grammarlens}[title={What you've built}]
One more word for saying what happens if.
\end{grammarlens}

\subsection*{Your turn}
\begin{itemize}
\raggedright
  \item \emph{Say it:} \emph{ādarśaḥ}
  \item \emph{Say it:} \emph{ādarśaḥ}, once more
  \item \emph{Say it:} \emph{yadi pitā satyavādī, tarhi saḥ putrasya ādarśaḥ bhavati.} — If a father is truthful, he is a model for his son.
  \item \emph{Recall:} say the Sanskrit for mild-mannered, then the Sanskrit for in vain, then say \emph{ādarśaḥ} again
\end{itemize}

\begin{tcolorbox}[breakable,colback=teal!4,colframe=teal!35!black,title={Before you move on}]
What does \sk{आदर्शः} mean? (\textquotedblleft{}An ideal; a model\textquotedblright{}.) Say it once more.
\end{tcolorbox}
